\documentclass[preprint,5p,times,twocolumn]{elsarticle}

\usepackage[utf8]{inputenc}
\usepackage[T1]{fontenc}
\usepackage{lmodern}
\usepackage{amsmath,amssymb,amsthm,mathtools}
\usepackage{amstext}
\usepackage{hyperref}
\usepackage[nameinlink]{cleveref}
\usepackage{enumitem}
\usepackage{booktabs}
\usepackage{array}
\usepackage{microtype}

\journal{Journal of Number Theory}

\hypersetup{
    colorlinks=true,
    linkcolor=blue,
    citecolor=blue,
    urlcolor=blue
}

\pdfstringdefDisableCommands{%
  \def\F{F}\def\Z{Z}\def\Q{Q}\def\R{R}\def\N{N}%
  \def\QR{QR}\def\NQR{NQR}}

\newtheorem{theorem}{Theorem}[section]
\newtheorem{lemma}[theorem]{Lemma}
\newtheorem{proposition}[theorem]{Proposition}
\newtheorem{remark}[theorem]{Remark}
\newtheorem{conjecture}[theorem]{Conjecture}
\newtheorem{definition}[theorem]{Definition}

\renewcommand{\qedsymbol}{$\blacksquare$}

\DeclareMathOperator{\ord}{ord}

\newcommand{\Z}{\mathbb{Z}}
\newcommand{\Q}{\mathbb{Q}}
\newcommand{\R}{\mathbb{R}}
\newcommand{\F}{\mathbb{F}}
\newcommand{\N}{\mathbb{N}}
\newcommand{\Leg}[2]{\left(\frac{#1}{#2}\right)}
\newcommand{\eps}{\varepsilon}
\newcommand{\QR}{\mathrm{QR}}
\newcommand{\NQR}{\mathrm{NQR}}

\bibliographystyle{elsarticle-num}

\begin{document}

\begin{frontmatter}

\title{An Explicit Evaluation of a Fibonacci Character Sum 
for Primes of Full Rank of Apparition}

\author[aff1]{Majid Ghandali}
\address[aff1]{Independent Researcher}

\begin{abstract}
Let $p\ge 7$ be a prime, $\chi_p$ the quadratic character mod $p$, 
and $(F_n)$ the Fibonacci sequence. The rank of apparition 
$\alpha(p)$ is the smallest $m\ge 1$ with $p\mid F_m$. 
Classically $\alpha(p)\mid p-\chi_p(5)$; hence the extremal 
condition $\alpha(p)=p-1$ forces $p\equiv\pm 1\pmod 5$.

In this full-rank regime we show first that $p\equiv 3\pmod 4$ 
and that the quadratic nonresidue root of $x^2-x-1$ in $\F_p$ 
is a primitive root of $\F_p^\times$. The Fibonacci character 
sum then collapses to
\[
\sum_{n=1}^{p-1}\chi_p(F_n)=
\begin{cases}+1,& p\equiv 11\pmod{20},\\
              -1,& p\equiv 19\pmod{20}.\end{cases}
\]
The proof reduces the sum to a single character value 
$\chi_p(r_- - r_+)$, where $r_-$ and $r_+$ denote the 
nonresidue and residue roots of $x^2-x-1$, via a 
primitive-root substitution; this value is then identified 
through a discriminant criterion linking the roots of 
$x^2-x-1$ to fifth roots of unity in $\F_p$. The result may 
be viewed as a Fibonacci analogue of an Artin-type primitivity 
condition. The theoretical result is accompanied by an 
independent computational verification up to $2\times 10^6$.
\end{abstract}

\begin{keyword}
Fibonacci numbers \sep Lucas sequences \sep character sums 
\sep rank of apparition \sep primitive roots \sep fifth roots of unity
\MSC[2020] 11B39 \sep 11A15 \sep 11L40 \sep 11Y60
\end{keyword}

\end{frontmatter}

\section{Introduction}\label{sec:intro}

The Fibonacci sequence $(F_n)_{n\ge 0}$, $F_0=0$, $F_1=1$, 
$F_{n+1}=F_n+F_{n-1}$, is the prototypical non-degenerate 
binary linear recurrence. For a prime $p\ne 5$, its mod-$p$ 
behaviour is encoded by two classical invariants: the 
\emph{Pisano period} $\pi(p)$ and the \emph{rank of apparition}
\[
\alpha(p):=\min\{m\ge 1:p\mid F_m\}.
\]
A classical result going back to Lucas \cite{Lucas1878} states
\begin{equation}\label{eq:rank-divides}
\alpha(p)\mid p-\chi_p(5),
\end{equation}
where $\chi_p$ denotes the quadratic character. Hence in the 
split case $\chi_p(5)=1$, equivalently $p\equiv\pm 1\pmod 5$, 
the maximal possible value of $\alpha(p)$ is $p-1$.
Throughout this paper, the term full-rank regime refers to the
extremal split value \(\alpha(p)=p-1\); the hypothesis itself will
force the split case.

\subsection{The Fibonacci character sum}

We study
\[
S(p):=\sum_{n=1}^{p-1}\chi_p(F_n).
\]
A priori $S(p)$ might be expected to fluctuate as $p$ varies, 
much like character sums attached to general linear recurrences 
\cite{Williams1982,Sun1992,Sun2003}. Our main theorem shows that 
in the extremal regime $\alpha(p)=p-1$ this apparent complexity 
collapses completely.

\begin{theorem}[Main Theorem]\label{thm:main}
Let $p\ge 7$ be a prime with $\alpha(p)=p-1$. Then 
$p\equiv 11\pmod{20}$ or $p\equiv 19\pmod{20}$, and
\[
S(p)=
\begin{cases}+1,& p\equiv 11\pmod{20},\\
              -1,& p\equiv 19\pmod{20}.\end{cases}
\]
\end{theorem}

\subsection{Why this regime?}

The hypothesis $\alpha(p)=p-1$ has an intrinsic primitive-root
interpretation. In the split case, \Cref{lem:rank-order} identifies
$\alpha(p)$ with the multiplicative order of $-r^2$, where $r$ is
either root of the Fibonacci polynomial
\[
x^2-x-1.
\]
Consequently, the full-rank condition is equivalent to the statement
that $-r^2$ is a primitive element of $\F_p^\times$. In this sense,
the problem belongs naturally to the same general circle of ideas as
Artin-type primitive-root questions.

The key structural fact proved in this paper is that the extremal
condition $\alpha(p)=p-1$ imposes significantly more structure than
the primitivity of $-r^2$ alone. We show that full rank forces the
split case, forces the congruence $p\equiv3\pmod4$, and implies that
the quadratic nonresidue root $r_-$ itself is primitive in
$\F_p^\times$. The mechanism is purely group-theoretic: since
$r_-$ is a quadratic nonresidue, one has
\[
r_-^{(p-1)/2}=-1,
\]
so $-1\in\langle r_-\rangle$. Since $r_-^2\in\langle r_-\rangle$,
it follows that
\[
-r_-^2\in\langle r_-\rangle.
\]
Because $-r_-^2$ already has order $p-1$, the subgroup generated by
$r_-$ must coincide with the entire multiplicative group
$\F_p^\times$. Thus $r_-$ is itself a primitive root.

This primitive-root phenomenon is the central structural input of the
paper. Once $r_-$ is known to be primitive, the sequence
\[
r_-^1,r_-^2,\dots,r_-^{p-1}
\]
runs through all elements of $\F_p^\times$ exactly once. The original
Fibonacci character sum
\[
S(p)=\sum_{n=1}^{p-1}\chi_p(F_n)
\]
can therefore be transformed from a character sum along a linear
recurrence sequence into a complete group sum over
$\F_p^\times$. This reindexing causes the entire expression to
collapse to the structural identity
\[
S(p)=\chi_p(r_- - r_+),
\]
where $r_+$ and $r_-$ denote the residue and nonresidue roots of
$x^2-x-1$.

The novelty of the present work is not merely the derivation of the
explicit formula of \Cref{thm:main}, but the discovery of this
structural mechanism linking three a priori distinct objects:
the rank of apparition of Fibonacci numbers, primitive roots in
finite fields, and quadratic character sums. The final evaluation of
$S(p)$ emerges as a consequence of this interaction rather than from a
direct character-sum computation. To the author's knowledge, neither
the primitive-root mechanism underlying the identity
\[
S(p)=\chi_p(r_- - r_+)
\]
nor the resulting explicit evaluation of $S(p)$ in the full-rank
regime has previously appeared in the literature.


\subsection{Outline of the proof}

The proof proceeds in three steps.
\begin{enumerate}[label=\textup{Step \Roman*.}, leftmargin=*]
\item \emph{Primitivity.}
We show that $\alpha(p)=p-1$ forces $x^2-x-1$ to split over 
$\F_p$, forces $p\equiv 3\pmod 4$, and makes the nonresidue root 
$r_-$ a primitive root of $\F_p^\times$ (\Cref{lem:G}).
\item \emph{Structural identity.}
Using Binet's formula and Step~I, the substitution $u=r_-^n$ 
converts the index $n$ into a complete traversal of $\F_p^\times$ 
and yields $S(p)=\chi_p(r_- - r_+)$ (\Cref{prop:structural}).
\item \emph{Explicit sign.}
A discriminant criterion linking the roots of $x^2-x-1$ to fifth 
roots of unity in $\F_p$ determines $\chi_p(r_- - r_+)$ in 
closed form (\Cref{lem:R3}).
\end{enumerate}

\subsection{Related literature}

The arithmetic of Fibonacci/Lucas sequences modulo primes begins 
with Lucas \cite{Lucas1878} and Wall \cite{Wall1960}; standard 
treatments are \cite{Vajda1989,Koshy2001,Lemmermeyer2000,
IrelandRosen1990}. Distribution of $\alpha(p)$ and related 
Pisano quantities is studied in \cite{Renault1996,BallotElia2007,
SomerKrizek2014,FiebigMbirikaSpilker2025}. Recent investigations 
on the $p$-adic and arithmetic structure of Lucas-type sequences 
include \cite{BragmanRowland2025,HajduTijdeman2025}. Character 
sums attached to linear recurrences appear in 
\cite{Williams1982,Sun1992,Sun2003,BerndtEvansWilliams1998}. 
The primitive-root perspective connects the full-rank condition 
with Artin-type problems; see Hooley's conditional density theorem 
\cite{Hooley1967} for the classical setting. The author is not aware 
of any previous explicit evaluation of $S(p)$ in the full-rank regime 
$\alpha(p)=p-1$.

\subsection{Organisation}

\Cref{sec:preliminaries} sets notation. \Cref{sec:structure} 
proves the primitivity result (Step~I). 
\Cref{sec:character-sums,sec:structural-identity} carry out 
Step~II. \Cref{sec:sign} executes Step~III. 
\Cref{sec:main-proof} assembles the main theorem. 
\Cref{sec:computations} reports the computational verification, 
and \Cref{sec:conclusion} states explicit density and 
generalisation questions.

\section{Preliminaries}\label{sec:preliminaries}

Throughout, $p$ denotes an odd prime $\ne 5$, with $p\ge 7$ 
unless explicitly stated otherwise. We use the standard extension 
of the quadratic character to all of $\F_p$ by setting 
$\chi_p(0)=0$, and recall that 
$\chi_p(-1)=(-1)^{(p-1)/2}$. We write
\[
\QR_p^\times=\{u\in\F_p^\times:\chi_p(u)=+1\},\qquad
\NQR_p=\{u\in\F_p^\times:\chi_p(u)=-1\}.
\]

When $p\equiv\pm 1\pmod 5$, the polynomial $x^2-x-1$ splits 
over $\F_p$. Fix a square root $s\in\F_p$ of $5$, and write 
its two roots as
\[
r_1=\frac{1+s}{2},\qquad r_2=\frac{1-s}{2}.
\]
Then $r_1+r_2=1$, $r_1\,r_2=-1$, and, with this choice of $s$, 
$r_1-r_2=s$. The Binet formula holds in $\F_p$:
\begin{equation}\label{eq:binet}
F_m=\frac{r_1^m-r_2^m}{r_1-r_2}.
\end{equation}
After the quadratic-character analysis in \Cref{lem:G}, in the 
full-rank case we shall denote the nonresidue root by $r_-$ and 
the residue root by $r_+$, so that
\[
\chi_p(r_-)= -1,\qquad \chi_p(r_+)=+1,
\]
and
\[
r_-+r_+=1,\qquad r_-\,r_+=-1.
\]
With this relabelling, the Binet formula may equivalently be 
written as
\[
F_m=\frac{r_-^m-r_+^m}{r_- - r_+}.
\]
The choice of $s$ will additionally be fixed by $\chi_p(s)=+1$ 
in \Cref{sec:sign}; this choice is independent of the labels 
$r_-$ and $r_+$, which are determined by quadratic character.

\subsection*{Elementary group-theoretic facts}

We shall use the following elementary facts about cyclic groups and
fifth roots of unity.

\begin{lemma}[Quadratic character in a cyclic group]
\label{lem:cyclic-character}
Let $g$ be a primitive root of $\F_p^\times$. Then, for every 
integer $k$,
\[
\chi_p(g^k)=(-1)^k.
\]
Equivalently, the even powers of $g$ are precisely the nonzero 
quadratic residues, and the odd powers of $g$ are precisely the 
quadratic nonresidues.
\end{lemma}

\begin{proof}
Since $g$ generates the cyclic group $\F_p^\times$ of order $p-1$, 
every element has the form $g^k$. The nonzero squares are exactly
\[
(g^j)^2=g^{2j}\qquad (j\in\Z),
\]
that is, precisely the even powers of $g$. Since $p-1$ is even, 
the parity of the exponent is well-defined modulo $p-1$. Hence 
$g^k$ is a square if and only if $k$ is even. Therefore 
$\chi_p(g^k)=+1$ for even $k$ and $\chi_p(g^k)=-1$ for odd $k$.
\end{proof}

\begin{lemma}[Primitive-root reindexing]
\label{lem:primitive-reindexing}
Let $g$ be a primitive root of $\F_p^\times$. Then the map
\[
\{1,2,\dots,p-1\}\longrightarrow \F_p^\times,\qquad n\mapsto g^n
\]
is a bijection. Consequently, for any function 
$H:\F_p^\times\to A$ with values in an abelian group $A$,
\[
\sum_{n=1}^{p-1} H(g^n)=\sum_{u\in\F_p^\times} H(u).
\]
Moreover,
\[
\chi_p(g^n)=(-1)^n\qquad(1\le n\le p-1).
\]
\end{lemma}

\begin{proof}
Since $g$ has order $p-1$, its powers $g,g^2,\dots,g^{p-1}$ are 
exactly the elements of $\F_p^\times$, each occurring once. This 
gives the asserted bijection and hence the reindexing identity. 
The final character formula follows from \Cref{lem:cyclic-character}.
\end{proof}

\begin{lemma}[Fifth roots and the cyclotomic polynomial]
\label{lem:phi5-order}
Let $K$ be a field of characteristic different from $5$, and let
\[
\Phi_5(X)=X^4+X^3+X^2+X+1.
\]
For $z\in K$, the following are equivalent:
\[
\Phi_5(z)=0
\qquad\Longleftrightarrow\qquad
z\in K^\times,\ z^5=1,\ z\ne 1.
\]
In particular, if $\Phi_5(z)=0$, then $z$ has multiplicative order
$5$.
\end{lemma}

\begin{proof}
The identity
\[
X^5-1=(X-1)\Phi_5(X)
\]
holds over every field. If $\Phi_5(z)=0$, then 
\[
z^5-1=(z-1)\Phi_5(z)=0,
\]
so $z^5=1$. Also
\[
\Phi_5(1)=5\ne0
\]
because the characteristic is not $5$; hence $z\ne1$. Thus $z$ is 
a nontrivial fifth root of unity. Its multiplicative order divides 
$5$ and is not $1$, so it is exactly $5$.

Conversely, if $z^5=1$ and $z\ne1$, then from 
\[
0=z^5-1=(z-1)\Phi_5(z)
\]
and $z-1\ne0$ we obtain $\Phi_5(z)=0$.
\end{proof}

\section{The structure lemmas (Step~I)}\label{sec:structure}

Throughout this section, $p$ denotes an odd prime with $p\ne5$. 
Whenever roots are used as elements of $\F_p$, the splitting 
hypothesis will be stated explicitly.

\begin{lemma}[Rank--order identity]\label{lem:rank-order}
Let $r_1,r_2$ be the two roots of $x^2-x-1$ in $\F_{p^2}$, and put
\[
q=\frac{r_1}{r_2}.
\]
Then
\[
\alpha(p)=\ord_{\F_{p^2}^\times}(q).
\]
In particular, if $x^2-x-1$ splits over $\F_p$, then for each root
$r\in\F_p$,
\[
\alpha(p)=\ord_{\F_p^\times}(-r^2).
\]
\end{lemma}

\begin{proof}
Since $p\ne5$, the roots are distinct. Binet's formula in 
$\F_{p^2}$ gives
\[
F_m=\frac{r_1^m-r_2^m}{r_1-r_2}.
\]
Hence
\[
F_m\equiv0\pmod p
\iff r_1^m=r_2^m
\iff \left(\frac{r_1}{r_2}\right)^m=1
\iff q^m=1.
\]
Therefore $\alpha(p)=\ord(q)$.

If the polynomial splits over $\F_p$ and $r$ is one of its roots, 
then the other root is $-r^{-1}$, since the product of the roots 
is $-1$. Thus the ratio of the two roots is either
\[
\frac{r}{-r^{-1}}=-r^2
\]
or its inverse. Therefore the order of this ratio is
$\ord_{\F_p^\times}(-r^2)$, proving the asserted split-case identity.
\end{proof}

\begin{lemma}[Primitive root mechanism in the full-rank regime]
\label{lem:G}
Let $p\ge 7$ be a prime such that $\alpha(p)=p-1$. Then
$x^2-x-1$ splits over $\F_p$, and $p\equiv 3\pmod 4$. Label its two
roots $r_+,r_-\in\F_p^\times$ by their quadratic characters:
\[
\chi_p(r_+)=+1,\qquad \chi_p(r_-)=-1.
\]
Then
\[
\ord_{\F_p^\times}(-r_-^2)=p-1
\]
and
\[
\ord_{\F_p^\times}(r_-)=p-1.
\]
In particular, $r_-$ is a primitive root of $\F_p^\times$.
\end{lemma}

\begin{proof}
Let $r_1,r_2$ be the two roots of $x^2-x-1$ in $\F_{p^2}$, and put
\[
q=\frac{r_1}{r_2}.
\]
By \Cref{lem:rank-order},
\[
\alpha(p)=\ord_{\F_{p^2}^\times}(q).
\]

Assume first that $x^2-x-1$ does not split over $\F_p$. Then Frobenius
interchanges the two roots, so
\[
q^p=\frac{r_1^p}{r_2^p}=\frac{r_2}{r_1}=q^{-1}.
\]
Hence $q^{p+1}=1$, and therefore $\ord(q)\mid p+1$. By
\Cref{lem:rank-order}, this gives
\[
\alpha(p)\mid p+1.
\]
Since $\alpha(p)=p-1$, we would have $p-1\mid p+1$, hence
$p-1\mid2$, impossible for $p\ge7$. Therefore $x^2-x-1$ splits over
$\F_p$.

Now let $r$ be either root in $\F_p$. By the split-case part of
\Cref{lem:rank-order},
\[
\ord_{\F_p^\times}(-r^2)=\alpha(p)=p-1.
\]
If $-1$ were a square modulo $p$, then $-r^2$ would be a square in
$\F_p^\times$. But the nonzero squares form a subgroup of
$\F_p^\times$ of order $(p-1)/2$, so no square can have order $p-1$.
This contradiction shows that
\[
\chi_p(-1)=-1,
\]
and hence $p\equiv3\pmod4$.

Since the product of the two roots is $-1$, we have
\[
\chi_p(r_1)\chi_p(r_2)=\chi_p(r_1r_2)=\chi_p(-1)=-1.
\]
Thus exactly one root is a quadratic residue and the other is a 
quadratic nonresidue. Label them as $r_+$ and $r_-$, respectively.

Applying the split-case identity in \Cref{lem:rank-order} to the root
$r_-$ gives
\[
\ord_{\F_p^\times}(-r_-^2)=p-1.
\]

It remains to show that $r_-$ itself is primitive. Since
$\chi_p(r_-)=-1$, Euler's criterion gives
\[
r_-^{(p-1)/2}=-1.
\]
Thus $-1\in\langle r_-\rangle$. Since $r_-^2\in\langle r_-\rangle$ 
as well, we have
\[
-r_-^2\in\langle r_-\rangle.
\]
But $-r_-^2$ has order $p-1$. Therefore the subgroup
$\langle r_-\rangle\subseteq\F_p^\times$ contains an element of order
$p-1$. Since $|\F_p^\times|=p-1$, it follows that
\[
\langle r_-\rangle=\F_p^\times.
\]
Hence
\[
\ord_{\F_p^\times}(r_-)=p-1,
\]
as claimed.
\end{proof}

\begin{remark}[The content of \Cref{lem:G}]\label{rem:content-G}
In the split case, the rank--order identity shows that the full-rank
condition is equivalent to the primitivity of $-r^2$ for either root
$r$ of $x^2-x-1$. The additional content of \Cref{lem:G} is twofold:
first, full rank itself forces the split case and the parity condition
$p\equiv3\pmod4$; second, for the nonresidue root $r_-$, the primitivity
of $-r_-^2$ forces the primitivity of $r_-$ itself. Indeed, since
$r_-$ is a nonresidue, $-1\in\langle r_-\rangle$, and hence
$-r_-^2\in\langle r_-\rangle$. This is the structural fact that makes the substitution
\(u=r_-^n\) in Step~II a bijection onto \(\F_p^\times\), and it is
the key structural input in the evaluation of \(S(p)\).
\end{remark}

\begin{remark}[Complementary order of the residue root]
Although not needed for the evaluation of $S(p)$, the order of the
residue root can also be read off in the full-rank regime. Write
$p=4m+3$ and $n=(p-1)/2=2m+1$. Since $r_-$ is primitive, 
$r_-^n=-1$. From $r_-r_+=-1$ we get
\[
r_+=-r_-^{-1}=r_-^n r_-^{-1}=r_-^{n-1}.
\]
Hence
\[
\ord_{\F_p^\times}(r_+)=\frac{2n}{\gcd(2n,n-1)}.
\]
Since $n-1=2m$,
\[
\gcd(2n,n-1)=\gcd(4m+2,2m)=2\gcd(2m+1,m)=2,
\]
and therefore
\[
\ord_{\F_p^\times}(r_+)=n=\frac{p-1}{2}.
\]
\end{remark}

\section{Character sum identities}\label{sec:character-sums}

\begin{lemma}\label{lem:S}
For every $a\in\F_p^\times$,
\[
\sum_{x\in\F_p}\chi_p(x^2-a)=-1.
\]
\end{lemma}

\begin{proof}
This is classical \cite[Ch.~6]{IrelandRosen1990}; we include a 
brief self-contained argument. Let 
\[
T(a):=\sum_{x\in\F_p}\chi_p(x^2-a).
\]
Using $\#\{x:x^2=t\}=1+\chi_p(t)$ and 
$\sum_t\chi_p(t-a)=0$, we obtain
\[
T(a)=\sum_{t\in\F_p}\bigl(1+\chi_p(t)\bigr)\chi_p(t-a)
=\sum_t\chi_p\bigl(t(t-a)\bigr).
\]
The substitution $t=aw$ in the last sum shows that, for 
$a\in\F_p^\times$, this value equals 
\[
\sum_{w\in\F_p}\chi_p(w(w-1)),
\]
and hence is independent of $a\ne0$. Since $T(0)=p-1$ and
\[
\sum_{a\in\F_p}T(a)=\sum_x\sum_a\chi_p(x^2-a)=0,
\]
the common value of $T(a)$ for $a\ne0$ must be $-1$.
\end{proof}

\begin{lemma}\label{lem:D}
If $p\equiv 3\pmod 4$, the map $u\mapsto -u$ is a bijection 
$\QR_p^\times\to\NQR_p$.
\end{lemma}

\begin{proof}
$\chi_p(-u)=\chi_p(-1)\chi_p(u)=-\chi_p(u)$, and the two sets 
have equal cardinality $(p-1)/2$.
\end{proof}

\begin{lemma}\label{lem:Q}
If $p\equiv 3\pmod 4$, then
\[
A:=\!\!\sum_{u\in\QR_p^\times}\!\!\chi_p(u^2-1)=0,\qquad
B:=\!\!\sum_{u\in\NQR_p}\!\!\chi_p(u^2+1)=-1.
\]
\end{lemma}

\begin{proof}
By \Cref{lem:S} with $a=1$, isolating the term $x=0$ 
(where $\chi_p(0^2-1)=\chi_p(-1)=-1$ since $p\equiv 3\pmod 4$):
\[
\sum_{x\in\F_p^\times}\chi_p(x^2-1)
=-1-\chi_p(-1)=-1-(-1)=0.
\]
The involution $u\mapsto -u$ on $\F_p^\times$ leaves $u^2-1$ 
invariant, and by \Cref{lem:D} swaps $\QR_p^\times$ and 
$\NQR_p$ bijectively. Hence the QR-half and the NQR-half of 
the sum are equal, and $2A=0$.

Similarly with $a=-1$, isolating $x=0$ 
(where $\chi_p(0^2+1)=\chi_p(1)=+1$):
\[
\sum_{x\in\F_p^\times}\chi_p(x^2+1)
=-1-\chi_p(1)=-1-1=-2.
\]
Since $(-u)^2+1=u^2+1$, the same involution argument gives 
$2B=-2$, hence $B=-1$.
\end{proof}

\section{The structural identity (Step~II)}\label{sec:structural-identity}

\begin{proposition}\label{prop:structural}
Let $p\ge 7$ with $\alpha(p)=p-1$. Then
\[
S(p)=\chi_p(r_- - r_+),
\]
where $r_-$ and $r_+$ are the nonresidue and residue roots of 
$x^2-x-1$ in $\F_p$.
\end{proposition}

\begin{proof}
By \Cref{lem:G}, the polynomial $x^2-x-1$ splits over $\F_p$, and
$r_-$ and $r_+$ are well-defined. Let
\[
D=r_- - r_+.
\]
Since the discriminant of $x^2-x-1$ is $5\ne 0$ in $\F_p$ 
(as $p\ne 5$), we have $D\ne 0$. By Binet, since $r_+=-r_-^{-1}$,
\[
F_n=\frac{r_-^n-(-1)^n r_-^{-n}}{D}.
\]
By \Cref{lem:G}, $r_-$ is a primitive root of $\F_p^\times$. 
Therefore, by \Cref{lem:primitive-reindexing}, the map 
\[
n\mapsto u=r_-^n
\]
sends $1,\dots,p-1$ bijectively onto $\F_p^\times$, and
\[
\chi_p(u)=\chi_p(r_-^n)=(-1)^n.
\]
Thus
\[
F_n=\frac{u-\chi_p(u)\,u^{-1}}{D},\qquad
\chi_p(F_n)=\eps\,\chi_p\!\bigl(u-\chi_p(u)u^{-1}\bigr),
\]
where $\eps:=\chi_p(D)$. We may therefore reindex the sum over 
$n$ as a sum over $u\in\F_p^\times$.

Splitting by quadratic character:
\begin{itemize}[leftmargin=*]
\item If $u\in\QR_p^\times$, then $\chi_p(u)=+1$ and 
\[
u-u^{-1}=\frac{u^2-1}{u},
\]
so the corresponding contribution is 
\[
\eps\,\chi_p(u^2-1).
\]

\item If $u\in\NQR_p$, then $\chi_p(u)=-1$ and 
\[
u+u^{-1}=\frac{u^2+1}{u},
\]
so the corresponding contribution is 
\[
-\eps\,\chi_p(u^2+1),
\]
because $\chi_p(u)=-1$.
\end{itemize}
By \Cref{lem:Q},
\[
S(p)=\eps(A-B)=\eps(0-(-1))=\eps=\chi_p(r_- - r_+).
\]
\end{proof}

\begin{remark}
The zero term at $n=p-1$ (where $u=1$, $u^2-1=0$) is absorbed 
automatically into the QR-sum and contributes $0$ to $A$.
\end{remark}

\section{Explicit sign via fifth roots of unity (Step~III)}
\label{sec:sign}

Throughout this section, assume $p\equiv 3\pmod 4$ and 
$p\equiv\pm 1\pmod 5$. Since $\chi_p(5)=1$, choose $s\in\F_p$ 
with $s^2=5$ and $\chi_p(s)=+1$ (the two square roots of $5$ 
have opposite characters because $\chi_p(-1)=-1$). Put
\[
\phi=\tfrac{1+s}{2},\qquad \psi=\tfrac{1-s}{2}.
\]
These are the two roots of $x^2-x-1$. At this stage they are not
labelled by quadratic character; the labels $r_-$ and $r_+$ will be
identified in \Cref{lem:R3}.

\begin{lemma}\label{lem:R0}
There exists $z\in\F_p^\times$ with $z^5=1$ and $z\ne 1$ if 
and only if $p\equiv 1\pmod 5$.
\end{lemma}

\begin{proof}
The group $\F_p^\times$ is cyclic of order $p-1$. It contains an
element of order $5$ if and only if $5\mid p-1$, equivalently
$p\equiv1\pmod5$. Such an element is precisely a nontrivial fifth
root of unity.
\end{proof}

\begin{lemma}[Discriminant criterion]\label{lem:R2}
Let $t=(s-1)/2$, and let
\[
\Delta=t^2-4
\]
be the discriminant of $Z^2-tZ+1$. Then $\Delta\ne0$ and
\[
\chi_p(\Delta)=
\begin{cases}+1,& p\equiv 1\pmod 5,\\
              -1,& p\equiv -1\pmod 5.\end{cases}
\]
\end{lemma}

\begin{proof}
The point is to translate the quadratic solvability of 
$z^2-tz+1$ into the existence of a nontrivial fifth root of unity 
in $\F_p$.

\emph{Setup.} From $s^2=5$, 
\[
t^2=\frac{(s-1)^2}{4}=\frac{3-s}{2},
\]
so
\begin{equation}\label{eq:t-relation}
t^2+t-1=0,\qquad \Delta=t^2-4=-t-3.
\end{equation}
If $\Delta=0$, then $t=-3$, and substituting into 
\eqref{eq:t-relation} gives $9-3-1=5=0$ in $\F_p$, contradicting 
$p\ne5$.

\smallskip
\emph{Step 1: Equivalence with a quadratic.} By the discriminant 
criterion for quadratics over $\F_p$,
\begin{equation}\label{eq:Delta-equiv}
\chi_p(\Delta)=+1 \;\iff\; \exists\,z\in\F_p:\; z^2-tz+1=0.
\end{equation}
Any such root $z$ lies in $\F_p^\times$ and satisfies 
$z+z^{-1}=t$.

\smallskip
\emph{Step 2: From $z+z^{-1}=t$ to an element of order $5$.}
Suppose $z\in\F_p^\times$ with $z+z^{-1}=t$. Then 
$z^2+z^{-2}=t^2-2$, and hence
\[
(z^2+z^{-2})+(z+z^{-1})+1
=(t^2-2)+t+1=t^2+t-1=0
\]
by \eqref{eq:t-relation}. Multiplying through by $z^2$ gives
\begin{equation}\label{eq:phi5}
z^4+z^3+z^2+z+1=\Phi_5(z)=0.
\end{equation}
By \Cref{lem:phi5-order}, $z$ has multiplicative order $5$. 
Therefore $\F_p^\times$ contains an element of order $5$, and by 
\Cref{lem:R0} we have $p\equiv1\pmod5$.

\smallskip
\emph{Step 3: Converse.} Assume $p\equiv1\pmod5$. By 
\Cref{lem:R0}, choose $z_0\in\F_p^\times$ of order $5$, and set 
\[
\tau:=z_0+z_0^{-1}.
\]
From $\Phi_5(z_0)=0$ (by \Cref{lem:phi5-order}), dividing by 
$z_0^2$ gives
\[
(z_0^2+z_0^{-2})+(z_0+z_0^{-1})+1=0.
\]
Using $z_0^2+z_0^{-2}=\tau^2-2$, we get
\[
\tau^2+\tau-1=0.
\]
Since $s^2=5$, the two roots of $X^2+X-1$ in $\F_p$ are
\[
\frac{-1+s}{2}=t,\qquad \frac{-1-s}{2}=-1-t.
\]
Thus $\tau\in\{t,-1-t\}$.

If $\tau=t$, then $z_0+z_0^{-1}=t$, so 
$z_0^2-tz_0+1=0$. By \eqref{eq:Delta-equiv}, 
$\chi_p(\Delta)=+1$.

If $\tau=-1-t$, replace $z_0$ by $z_1:=z_0^2$, which still has 
order $5$. Then
\[
z_1+z_1^{-1}=z_0^2+z_0^{-2}=\tau^2-2.
\]
Since $\tau^2+\tau-1=0$, we have $\tau^2=1-\tau$. For 
$\tau=-1-t$, this gives
\[
\tau^2=1-(-1-t)=2+t,
\]
and hence
\[
z_1+z_1^{-1}=\tau^2-2=t.
\]
Therefore $z_1^2-tz_1+1=0$, and again by 
\eqref{eq:Delta-equiv} we obtain $\chi_p(\Delta)=+1$.

Combining Steps 2 and 3 gives 
\[
\chi_p(\Delta)=+1 \quad\Longleftrightarrow\quad p\equiv1\pmod5.
\]
Since $\Delta\ne0$ and throughout this section 
$p\equiv\pm1\pmod5$, the asserted two-case formula follows.
\end{proof}

\begin{lemma}[Sign determination]\label{lem:R3}
With $r_-$ the nonresidue root and $r_+$ the residue root of 
$x^2-x-1$,
\[
\chi_p(r_- - r_+)=
\begin{cases}+1,& p\equiv 1\pmod 5,\\
              -1,& p\equiv -1\pmod 5.\end{cases}
\]
\end{lemma}

\begin{proof}
Compute
\[
s\phi=s\frac{1+s}{2}=\frac{s+s^2}{2}=\frac{s+5}{2}.
\]
By \eqref{eq:t-relation},
\[
\Delta=-t-3=-\frac{s-1}{2}-3=-\frac{s+5}{2}=-s\phi.
\]
Taking characters and using $\chi_p(s)=+1$ and $\chi_p(-1)=-1$,
\[
\chi_p(\Delta)=\chi_p(-1)\chi_p(s)\chi_p(\phi)=-\chi_p(\phi).
\]

\emph{Case $p\equiv 1\pmod 5$.} By \Cref{lem:R2}, 
$\chi_p(\Delta)=+1$, so $\chi_p(\phi)=-1$. Hence $r_-=\phi$, 
$r_+=\psi$, and
\[
r_- - r_+ = \phi-\psi = s.
\]
Therefore
\[
\chi_p(r_- - r_+) = \chi_p(s) = +1.
\]

\emph{Case $p\equiv -1\pmod 5$.} By \Cref{lem:R2}, 
$\chi_p(\Delta)=-1$, so $\chi_p(\phi)=+1$. Hence $r_+=\phi$, 
$r_-=\psi$, and
\[
r_- - r_+ = \psi-\phi = -s.
\]
Therefore
\[
\chi_p(r_- - r_+) = \chi_p(-1)\chi_p(s) = -1.
\]
\end{proof}

\section{Proof of the main theorem}\label{sec:main-proof}

\begin{proof}[Proof of \Cref{thm:main}]
By \Cref{lem:G}, the polynomial $x^2-x-1$ splits over $\F_p$ and
$p\equiv3\pmod4$. Since the discriminant of $x^2-x-1$ is $5$, the
splitting condition is equivalent to $\chi_p(5)=+1$, hence
$p\equiv\pm1\pmod5$. Combining these congruences gives
\[
p\equiv 11\pmod{20}\quad\text{or}\quad p\equiv 19\pmod{20}.
\]

By \Cref{prop:structural}, 
\[
S(p)=\chi_p(r_- - r_+).
\]
By \Cref{lem:R3}, this equals $+1$ when $p\equiv 1\pmod 5$ 
(i.e.\ $p\equiv 11\pmod{20}$) and $-1$ when $p\equiv -1\pmod 5$ 
(i.e.\ $p\equiv 19\pmod{20}$).
\end{proof}

\section{Computational verification}\label{sec:computations}

\Cref{thm:main} was verified for all primes \(p\le 2\times 10^6\)
satisfying \(\alpha(p)=p-1\).


\begin{table}[ht]\centering
\caption{Verification of \Cref{thm:main}.}\label{tab:verification}
\begin{tabular}{ccccc}\toprule
$p\bmod 20$ & predicted $S(p)$ & count & matches & status\\
\midrule
$11$ & $+1$ & $11{,}755$ & $11{,}755$ & PASS\\
$19$ & $-1$ & $14{,}652$ & $14{,}652$ & PASS\\
\midrule
total &     & $26{,}407$ & $26{,}407$ & PASS\\
\bottomrule
\end{tabular}\end{table}

The verification is non-circular: the program computes 
$\alpha(p)$ by iterating the Fibonacci recurrence modulo $p$ 
until the first zero is reached, and computes $S(p)$ by 
directly summing $\chi_p(F_n)$ for $1\le n\le p-1$, without 
using the closed formula of \Cref{thm:main}. Several internal 
consistency checks were also performed:
\begin{itemize}[leftmargin=*,nosep]
\item[(i)] the quadratic characters of the two roots of $x^2-x-1$ 
in $\F_p$ were verified against \Cref{lem:G};
\item[(ii)] the multiplicative orders of $-r_-^2$ and $r_-$ were 
checked against \Cref{lem:rank-order,lem:G};
\item[(iii)] the sign of $\chi_p(r_- - r_+)$ was 
cross-checked against the discriminant criterion of 
\Cref{lem:R3}.
\end{itemize}

These checks verify not only the numerical value of the final sum, 
but also the structural mechanism used in the proof: the passage from 
full rank to the primitivity of $r_-$, the structural identity 
$S(p)=\chi_p(r_- - r_+)$, and the final sign determination.

Among the $148{,}933$ primes below $2\times 10^6$, exactly 
$26{,}407$ satisfy $\alpha(p)=p-1$. Restricting to the split 
class $p\equiv\pm 1\pmod 5$, which contains $74{,}461$ primes 
in this range, the full-rank density is 
$26{,}407/74{,}461\approx 35.46\%$. Source code, datasets and 
verification logs are deposited in the reproducibility archive 
\cite{Ghandali2026}.

\section{Concluding remarks}\label{sec:conclusion}

We have established an explicit evaluation of the Fibonacci 
character sum $S(p)$ in the full-rank regime $\alpha(p)=p-1$, 
showing that $S(p)=\pm 1$ with sign determined by $p\bmod 20$. 
Since full rank forces the split condition $\chi_p(5)=1$, the
split-case identity in \Cref{lem:rank-order} identifies
$\alpha(p)=p-1$ with the primitivity of $-r^2\in\F_p^\times$ for
either root $r$ of $x^2-x-1$. This places the result naturally in
the Artin-type family.

\subsection*{A density question}

Among primes \(p\le 2\times 10^6\) in the split congruence classes
\(p\equiv\pm1\pmod 5\) (the only mod-\(5\) classes in which
\(\alpha(p)=p-1\) can occur), our data show that approximately
\(35.46\%\) satisfy the full-rank condition.\Cref{thm:main} restricts these primes further to 
the two classes $p\equiv 11,19\pmod{20}$, among the four split 
classes $p\equiv 1,9,11,19\pmod{20}$. This structural constraint 
motivates the following.

\begin{conjecture}\label{conj:density}
There exists a constant $C_F$ with $0<C_F\le 1/2$ such that
\[
\#\{p\le x:\alpha(p)=p-1\}\sim 
C_F\cdot\#\{p\le x:p\equiv\pm 1\pmod 5\},\quad x\to\infty.
\]
Equivalently, if this asymptotic holds, the natural density of 
full-rank primes among all primes equals $C_F/2$.
\end{conjecture}

The upper bound $C_F\le 1/2$ is forced by \Cref{thm:main}: 
full-rank primes lie in the two classes $\{11,19\}\pmod{20}$ 
out of the four split classes $\{1,9,11,19\}\pmod{20}$.By the split-case identity in \Cref{lem:rank-order}, the conjecture
is a primitive-root statement for the reductions, at primes split in
\(\mathbb Q(\sqrt5)\), of the algebraic unit
\(-\phi^2\), where \(\phi=(1+\sqrt5)/2\), equivalently of its
conjugate. A Hooley-type 
analysis in this conditional setting, under suitable GRH-type 
hypotheses \cite{Hooley1967}, would be expected to yield a positive 
value of $C_F$ given by an Euler product encoding the splitting 
behaviour of $-r^2$ in the relevant Galois extensions. A rigorous 
treatment lies beyond the scope of this paper.


\subsection*{Further directions}

The methods of this paper apply directly to related 
Lucas-type character sums via the substitution 
$L_n=r_-^n+r_+^n$; a detailed development of such analogues, 
together with the analysis of $S(p)$ in non-extremal regimes, 
is left to future work.

\appendix

\section{Illustrative examples}\label{app:examples}

For $p=11$: $\sqrt 5\equiv 4\pmod{11}$, so the roots of 
$x^2-x-1$ are $8$ and $4$; here $r_-=8$ (nonresidue), 
$r_+=4$ (residue), $r_- - r_+ = 4 = 2^2$, and $S(11)=+1$.\\
For $p=19$: roots $15$ and $5$; $r_-=15$, $r_+=5$, 
$r_- - r_+ = 10$, a nonresidue, and $S(19)=-1$.\\
For $p=31$: $\alpha(31)=30$, and the roots are $13$ and $19$; 
here $r_-=13$, $r_+=19$, $r_- - r_+\equiv 25=5^2\pmod{31}$, 
and $S(31)=+1$.\\
Direct computation also gives $\alpha(59)=58$ with $S(59)=-1$, 
and $\alpha(79)=78$ with $S(79)=-1$, in agreement with 
\Cref{thm:main}.

\section{Reproducibility archive}\label{app:software}

Complete source code, datasets, and verification logs are 
deposited at \cite{Ghandali2026}.

\bibliography{references}

\end{document}
